% ============================================================================
% 02_related.tex — Related Work (merged, six blocks per plan.md §3)
% Sources: writeup/contents/related_work.tex (blocks i, iii, v, vi)
%          Engineering_simplicity/main.tex:93–112 (blocks ii, iii, iv, v)
% ============================================================================

\section{Related Work}\label{sec:related}

\paragraph{Auction theory and experiments.}
There is a vast quantity of theoretical and experimental literature on auctions. While \citet{krishna2009auction}'s textbook and \citet{kagel2020handbook}'s handbook provide invaluable general resources, we focus on citations relevant to the results in this paper.
Starting with the IPV case, the benchmark of revenue equivalence in the risk-neutral case is developed in \citet{myerson1981optimal}. The experimental evidence for departures due to risk-aversion in the FPSB auction is documented in \citet{coppinger1980incentives} and \citet{cox1988theory}'s survey. The experimental evidence for the common error of bidding \emph{above} one's value in the SPSB auction is documented in \citet{kagel1993independent}; as we show in Section~\ref{sec:validation}, LLM bidders instead tend to bid \emph{below} value in the SPSB auction---a qualitative divergence from human behavior whose implications for the design ranking we take up in Section~\ref{sec:humans}.
For common value settings, experimental evidence for the winner's curse was first given by \citet{bazerman1983won}; we primarily follow the experimental design and analysis framework in \citet{kagel1986winner} and \citet{kagel1987information}. \citet{levin1996revenue} extend this line to English (ascending) common-value auctions, where the information released by observed dropouts attenuates, but does not eliminate, the curse. For eBay-style auctions, the phenomenon of bid sniping has been studied by \citet{roth2002last} and \citet{ockenfels2006late}, while \citet{einav2011learning} provide an empirical study of hidden reserves on eBay.

\paragraph{Simple mechanism design.}
\citet{li2017obviously} formalized obvious strategy-proofness (OSP) and showed that ascending clock auctions satisfy it while SPSB auctions do not. Related notions include \textit{strategic simplicity} \citep{borgers2019strategically}, which requires that optimal play not depend on higher-order beliefs, and \textit{one-step simplicity} \citep{pycia2023theory}, which requires that optimal play emerge from considering only the immediate next step rather than full backward induction. The SPSB auction satisfies strategic simplicity---truthful bidding is dominant regardless of beliefs about others---but is not OSP. Ascending clock auctions satisfy both. \citet{li2024designing} decomposes mechanism complexity into three cognitive dimensions: contingent reasoning, forward planning, and belief formation. We use this decomposition to structure the reasoning-scaffold family of our lever taxonomy (Section~\ref{sec:framework}).

Recent work emphasizes the connection between strategic complexity and apparent irrationality. \citet{oprea2024decisions} argues that many anomalies attributed to non-standard preferences are better explained by computational complexity: subjects make mistakes not because their preferences are irrational but because the problem is hard. This framing aligns with our approach. We do not assume LLMs have ``biases'' in the behavioral economics sense (though we probe prospect-theoretic preference framings in Appendix~\ref{app:prospect}); rather, our primary question is whether the computational structure of strategic problems creates failure modes for LLM agents---and which design levers remove them.

\paragraph{Mechanism descriptions and framing.}
A complementary literature holds the mechanism fixed and varies only how it is described. \citet{gonczarowski2023strategyproofness} study descriptions of strategy-proof mechanisms that aim to make incentive properties transparent, presenting short menu-based descriptions in which a participant's report is used only to select from a menu determined by others' reports; they develop such descriptions both for the SPSB auction and for deferred acceptance (DA), though their menu-description treatment of an SPSB auction did not perform well empirically. In an incentivized lab experiment, \citet{gonczarowski2024describing} find that a menu-based explanation significantly improved participants' measured understanding of DA's strategyproofness relative to standard descriptions, though average behavioral effects were modest. \citet{breitmoser2022obviousness} shows that re-framing a static sealed-bid auction as an ascending clock auction---a change of description, not of extensive form---can improve the rate of truth-telling behavior. Three further results discipline how description treatments should be read. \citet{guillen2018effectiveness} find in a field experiment on top trading cycles that describing the strategy-proofness \emph{property} increases truth-telling while describing the mechanism's \emph{mechanics} significantly reduces it; \citet{katuscak2024simpler} show that a description isolating what does and does not depend on a participant's own report increases truthful play in TTC, more so for higher-numeracy subjects; and \citet{masuda2022net} show that direct \emph{advice} to bid truthfully in a Vickrey auction raises truthful bidding by roughly 24 percentage points net of experimenter-demand effects---a treatment type we deliberately exclude from our lever space, since advice reveals optimal play rather than exposing why it is safe. These strands---menu restatements with weak behavioral effects, framings that import the salience of a dynamic format, and content that exposes the incentive-relevant invariant as opposed to the mechanics---correspond to distinct description sub-families in our taxonomy (Section~\ref{sec:framework}) and to distinct rungs of our ranking (Section~\ref{sec:ranking-descriptions}).

\paragraph{Matching and deferred acceptance.}
There is a deep literature in market design for two-sided matching. \citet{gale1962college} introduced the DA algorithm, which produces stable matchings, and \citet{roth1982economics} showed that DA is strategy-proof for the proposing side.
The question of simplicity in DA was addressed first by \citet{ashlagi2018stable}, who show that DA is generally not OSP-implementable. However, under \textit{acyclic} priority structures---a condition introduced by \citet{ergin2002efficient}---the proposer-optimal stable matching rule can be implemented in an OSP manner (for the proposer). Acyclicity rules out configurations where student $a$ has priority over $b$ at school $s$, while $b$ has priority over $c$ at school $s'$, and $c$ has priority over $a$ at school $s''$. Under acyclic priorities, rejection chains remain local, enabling an iterative query protocol that is OSP. Our iterative DA treatment implements exactly this protocol; to our knowledge, it constitutes the first experimental test of an OSP implementation of DA with any kind of agent. In this paper the matching domain serves as a robustness and generalization check on the auction results: the full lever grid is replicated in DA, with results reported in Appendix~\ref{app:da}.

\paragraph{LLMs as economic agents.}
While LLMs are able to generate human-like text and reasoning \citep{achiam2023gpt, bubeck2023sparks}, results also show that they can have limited planning abilities and exhibit various cognitive biases \citep{wan2023kelly}. A growing literature explores using these human-like AI models as simulated agents in economics and social science. \citet{horton2023large} introduces the ``homo silicus'' framework, showing that LLMs can replicate qualitative patterns from classic experiments in labor economics, bargaining, and social preferences; \citet{aher2023using} demonstrate that LLMs can reproduce human-subject results across a range of paradigms; \citet{brand2023using} use LLMs for market research, finding that they approximate human consumer behavior; and \citet{park2023generative} and \citet{manning2024automated} develop simulation and automated-experimentation infrastructure around LLM agents. Beyond simulation, there is growing interest in connecting LLMs with game-theoretic modeling of strategic environments more broadly---for example, deriving extensive-form structure, incentives, or interaction graphs from unstructured narratives \citep{daskalakis2024charting}, or integrating LLMs into mechanism-design pipelines and general-sum strategic settings \citep{huang2025accelerated, soumalias2025llm, jiang2025incentive, wang2026llm}.

Within auction settings specifically, there has been a proliferation of work since early versions of this paper, which can be broadly divided into three categories.

\emph{LLMs as first-class participants.} The first category treats LLMs as first-class entities---bidders or agents who actively participate in a marketplace---and studies their strategic behavior in its own right. \citet{fish2024algorithmic} study collusive behavior in first-price sealed-bid auctions. \citet{chen2023put} study how to make LLMs better at playing auctions. More recent work includes studies of information disclosure with LLM agents \citep{infobid2025}, comparisons across frontier models in various auction formats \citep{lu2025claude}, and early feasibility checks on LLM bidding behavior \citep{tsuchihashi2023homo}. \citet{auctionnet2024} also study a budget-constrained repeated ad auction (GSP-style), where bidders optimize campaign objectives over time. Much of this work uses auctions primarily as agent benchmarks for planning or collusion, without validating distributional similarity to human lab bidding, or focuses on narrow mechanism families without cross-format replication.

\emph{LLMs as proxies for human behavior in economic design.} The second category asks whether LLM agents can serve as useful proxies for human participants, with the goal of informing mechanism design. The closest work to ours here is \citet{corrigan2025usefulness}, who take distributional similarity seriously as the criterion of usefulness and show that a light-touch calibration step using a small sample of human bids can bring LLM bid distributions closer to human benchmarks. Their calibration agenda is complementary to ours and directly anchors the question we take up in Section~\ref{sec:humans}: our validation layer (Section~\ref{sec:validation}) benchmarks out-of-the-box behavior across multiple canonical formats against moment-matched human data---the ordinal fidelity evidence that licenses using LLM populations to rank design levers---and Section~\ref{sec:humans} then asks the \citeauthor{corrigan2025usefulness}-style question of what it takes to agree \emph{cardinally}, finding that the best tuning knobs are mechanism- and model-specific rather than portable.

\emph{Auction design for LLM-powered platforms.} A third, distinct line of work studies the design of auction mechanisms \emph{within} LLM-powered platforms---for example, how to allocate and price advertising slots or sponsored content when an LLM generates the display surface \citep{duetting2024mechanism, ramage2024adauctions}. Here the research question is about mechanism design for a new marketplace, rather than about how design choices shape the behavior of LLM participants, and is thus orthogonal to our agenda.

Our contribution differs in focus from all three categories: we use the conceptual vocabulary of simple mechanism design---contingent reasoning, forward planning, belief formation---to probe which implementation, description, and scaffolding levers preserve dominant-strategy play for LLM agents, and to ask whether the same theory that was built to characterize bounded rationality in humans also organizes theirs.

\paragraph{Automated and computational mechanism design.}
Our work also connects to computational and automated mechanism design, which uses algorithms and optimization to search over (or learn) allocation and payment rules in complex environments. A prominent learning-based line trains neural networks to maximize an objective such as revenue while controlling incentive-compatibility violations---as in RegretNet and its successors \citep{dutting2024optimal, curry102automated} and architectures such as GemNet that build strategyproofness into the mechanism class itself \citep{wang2024gemnet}---with extensions to auctions with budgets, combinatorial auctions, and data markets \citep{feng2018deep, wang2025bundleflow, ravindranath2023data}. Our contribution is orthogonal: rather than learning a new mechanism, we measure how implementation, description, and scaffolding choices affect the behavior of a fixed population of agents, and the resulting simplicity-preservation ranking can be used as an input to (and behavioral diagnostic for) automated design pipelines.

\bigskip

Relative to these literatures, our contribution is a single empirical object: a ranking of the simplicity-preserving levers available to a mechanism designer---extensive-form changes, mechanism descriptions, and reasoning scaffolds---measured on a common metric in canonical auction environments across four LLM families (Section~\ref{sec:ranking}) and replicated in a school-choice matching domain (Appendix~\ref{app:da}), licensed by benchmarking LLM behavior against moment-matched reconstructions of decades of human experimental evidence (Section~\ref{sec:validation}), and anchored rung-by-rung to human data wherever a human experiment exists (Section~\ref{sec:humans}).
